% !TEX root = ../main.tex
% ============================================================
% Chapter 5 -- The Frequency Transfer Operator
% Revised after discussion with Kees.
% Focus: conceptual and mathematical bridge from frequency coherence
% to residual-compatible learned solver components.
% ============================================================

\chapter{The frequency transfer operator}
\label{ch:frequency-transfer}

\providecommand{\chfiveincludegraphics}[2]{%
    \IfFileExists{#2}{%
        \includegraphics[width=#1]{#2}%
    }{%
        \fbox{%
            \parbox[c][0.22\textheight][c]{0.86\linewidth}{%
                \centering
                \textbf{Figure placeholder}\\[0.6em]
                \texttt{\detokenize{#2}}%
            }%
        }%
    }%
}

The preceding chapters introduced the Helmholtz equation, its finite
difference discretisation, absorbing boundary treatment, and the classical
iterative methods used later in the thesis. The central numerical difficulty
is that high-frequency Helmholtz systems are indefinite, oscillatory, and hard
to precondition uniformly. Classical methods act at one frequency at a time.
They do not exploit the fact that wavefields at neighbouring frequencies are
not unrelated objects.

This chapter introduces frequency transfer: the use of information from a
lower-frequency Helmholtz solve to assist a higher-frequency solve. The idea
is physically natural. In a homogeneous medium, changing the frequency changes
the wavelength and amplitude scale, but not the underlying propagation
geometry. The low-frequency field therefore contains information about where
energy travels, reflects, and decays. A learned operator can, in principle,
use this coherent structure to predict or correct the high-frequency field.

The main message of the chapter is deliberately more restrictive than this
physical motivation. Frequency transfer is plausible as a field-prediction
problem, but a Helmholtz solver is not controlled by field error alone. A
Krylov method is controlled by residuals. Therefore, the object needed by a
solver is not merely a map from one solution field to another. It is a
residual-compatible correction mechanism. This distinction is the backbone of
the thesis.

The chapter is organised as follows.
Section~\ref{sec:ch5-frequency-coherence} explains frequency coherence from
the Green's function. Section~\ref{sec:ch5-solution-transfer} defines exact
solution transfer and its Fourier interpretation. Section~\ref{sec:ch5-warm-starts}
shows how solution transfer can be used as a warm start, and why the induced
residual is the relevant algebraic test. Section~\ref{sec:ch5-transfer-precond}
then distinguishes solution transfer, residual restriction, and correction
prolongation. Section~\ref{sec:ch5-modal-safety} gives the modal mechanism by
which small field errors can become large residuals. Section~\ref{sec:ch5-pml}
explains why this issue is especially important in PML discretisations.
Finally, Sections~\ref{sec:ch5-objectives}--\ref{sec:ch5-architecture}
summarise the implications for learning objectives and neural architectures.

% ============================================================
\section{Frequency coherence of Helmholtz fields}
\label{sec:ch5-frequency-coherence}
% ============================================================

Consider the time-harmonic Helmholtz equation in a homogeneous medium with
unit wave speed,
\begin{equation}
    \mathcal{A}_\omega u
    :=
    \left(-\Delta-\omega^2\right)u
    =
    f,
    \label{eq:ch5-continuous-helmholtz}
\end{equation}
where $\omega$ is the angular frequency. In two-dimensional free space, the
outgoing Green's function for a point source at $\mathbf{x}_s$ is
\begin{equation}
    G_\omega(\mathbf{x},\mathbf{x}_s)
    =
    \frac{i}{4}H_0^{(1)}(\omega r),
    \qquad
    r=\|\mathbf{x}-\mathbf{x}_s\|_2,
    \label{eq:ch5-greens-function}
\end{equation}
where $H_0^{(1)}$ is the Hankel function of the first kind. For large
argument,
\begin{equation}
    H_0^{(1)}(z)
    \sim
    \sqrt{\frac{2}{\pi z}}
    \exp\!\left(i(z-\pi/4)\right),
    \qquad |z|\to\infty.
    \label{eq:ch5-hankel-asymptotic}
\end{equation}
Substituting~\eqref{eq:ch5-hankel-asymptotic} into
\eqref{eq:ch5-greens-function} gives the far-field form
\begin{equation}
    G_\omega(\mathbf{x},\mathbf{x}_s)
    \sim
    C_\omega r^{-1/2}e^{i\omega r},
    \qquad
    C_\omega
    =
    \frac{i}{4}
    \sqrt{\frac{2}{\pi\omega}}
    e^{-i\pi/4}.
    \label{eq:ch5-green-far-field}
\end{equation}

Equation~\eqref{eq:ch5-green-far-field} separates the field into three
components. The factor $r^{-1/2}$ is the two-dimensional geometric spreading
envelope. It depends on propagation distance, but not on the frequency. The
constant $C_\omega$ contains the global frequency-dependent amplitude scale,
with $|C_\omega|\propto\omega^{-1/2}$. The phase factor $e^{i\omega r}$
contains the oscillation rate. Increasing $\omega$ therefore increases the
number of phase cycles along the same ray path while preserving the large-scale
propagation geometry.

For a doubled-frequency pair $\omega_H=2\omega_L$, the far-field amplitude
ratio is
\begin{equation}
    \frac{|C_{\omega_H}|}{|C_{\omega_L}|}
    =
    \sqrt{\frac{\omega_L}{\omega_H}}
    =
    \frac{1}{\sqrt{2}}.
    \label{eq:ch5-amplitude-ratio}
\end{equation}
This ratio should not be overinterpreted. A high-frequency field is not a
pointwise rescaling of the low-frequency field, because the phase changes
along every propagation path. The important point is weaker and more useful:
the two solutions are coherent. They share geometry, but differ in phase,
amplitude scale, and boundary-layer behaviour.

\begin{figure}[htbp]
    \centering
    \chfiveincludegraphics{0.92\linewidth}{figures/ch5/green_frequency_coherence.png}
    \caption{Frequency coherence of homogeneous Helmholtz fields. The real
    part of the free-space Green's function is shown for several frequencies
    with one centred source. The propagation geometry is common to all panels,
    while the oscillation rate increases with frequency. Each panel should be
    normalised by its maximum absolute value so that the visual comparison
    emphasises phase structure rather than absolute amplitude.}
    \label{fig:ch5-green-frequency-coherence}
\end{figure}

\begin{figure}[htbp]
    \centering
    \chfiveincludegraphics{0.78\linewidth}{figures/ch5/radial_green_cross_section.png}
    \caption{Radial cross-section of the free-space Green's function for a
    doubled-frequency pair. The amplitude envelope should show the far-field
    scaling $\sqrt{\omega_L/\omega_H}$, while the real part should show the
    doubled oscillation rate. A shaded near-source region can mark where the
    far-field approximation is least accurate.}
    \label{fig:ch5-radial-green-cross-section}
\end{figure}

This physical coherence is the reason a data-driven transfer operator is a
reasonable object to learn. The low-frequency field is a compressed
description of propagation geometry. The high-frequency field asks for the
same geometry at a shorter wavelength.

% ============================================================
\section{Exact solution transfer}
\label{sec:ch5-solution-transfer}
% ============================================================

Let $u_L$ and $u_H$ solve two Helmholtz problems with the same source,
\begin{equation}
    \mathcal{A}_L u_L=f,
    \qquad
    \mathcal{A}_H u_H=f,
    \label{eq:ch5-two-frequency-solves}
\end{equation}
where $\mathcal{A}_L=\mathcal{A}_{\omega_L}$ and
$\mathcal{A}_H=\mathcal{A}_{\omega_H}$. If both operators are invertible under
the chosen boundary or radiation condition, the exact low-to-high solution
transfer operator is
\begin{equation}
    \mathcal{T}_{\uparrow}^{\mathrm{sol}}
    :=
    \mathcal{A}_H^{-1}\mathcal{A}_L.
    \label{eq:ch5-exact-transfer-up}
\end{equation}
Indeed,
\begin{equation}
    \mathcal{T}_{\uparrow}^{\mathrm{sol}}u_L
    =
    \mathcal{A}_H^{-1}\mathcal{A}_L u_L
    =
    \mathcal{A}_H^{-1}f
    =
    u_H.
    \label{eq:ch5-transfer-maps-solution}
\end{equation}
Similarly, the exact high-to-low solution-transfer operator is
\begin{equation}
    \mathcal{T}_{\downarrow}^{\mathrm{sol}}
    :=
    \mathcal{A}_L^{-1}\mathcal{A}_H.
    \label{eq:ch5-exact-transfer-down}
\end{equation}

The superscript ``sol'' is important. The operator in
\eqref{eq:ch5-exact-transfer-up} maps one complete solution field to another
complete solution field. It satisfies the intertwining identity
\begin{equation}
    \mathcal{A}_H\mathcal{T}_{\uparrow}^{\mathrm{sol}}
    =
    \mathcal{A}_L.
    \label{eq:ch5-solution-intertwining}
\end{equation}
On the solution manifold, this identity says that applying the high-frequency
operator after exact transfer recovers the original source $f$.

One should not replace~\eqref{eq:ch5-solution-intertwining} by an informal
commutation relation such as
$\mathcal{A}_H\mathcal{T}=\mathcal{T}\mathcal{A}_L$. Such a relation would
require a single operator to transfer both fields and right-hand sides in a
compatible way. That is not what a solution-transfer map is. This distinction
becomes decisive when the map is placed inside a Krylov method, because a
Krylov preconditioner acts on residuals rather than on complete solution
fields.

\begin{figure}[htbp]
    \centering
    \chfiveincludegraphics{0.86\linewidth}{figures/ch5/solution_transfer_operator_schematic.png}
    \caption{Exact solution transfer between two frequencies. The same source
    $f$ defines both $u_L=\mathcal{A}_L^{-1}f$ and
    $u_H=\mathcal{A}_H^{-1}f$. The operator
    $\mathcal{T}_{\uparrow}^{\mathrm{sol}}=\mathcal{A}_H^{-1}\mathcal{A}_L$
    maps $u_L$ to $u_H$ and satisfies
    $\mathcal{A}_H\mathcal{T}_{\uparrow}^{\mathrm{sol}}u_L=f$ on this
    solution pair.}
    \label{fig:ch5-solution-transfer-schematic}
\end{figure}

In a constant-coefficient periodic or free-space setting, the Fourier
transform diagonalises the Helmholtz operator:
\begin{equation}
    \widehat{\mathcal{A}_\omega u}(\mathbf{k})
    =
    \left(|\mathbf{k}|^2-\omega^2\right)\hat{u}(\mathbf{k}).
    \label{eq:ch5-fourier-symbol-A}
\end{equation}
The formal Fourier multiplier for exact solution transfer is therefore
\begin{equation}
    \widehat{\mathcal{T}_{\uparrow}^{\mathrm{sol}}}(\mathbf{k})
    =
    \frac{|\mathbf{k}|^2-\omega_L^2}
         {|\mathbf{k}|^2-\omega_H^2},
    \label{eq:ch5-fourier-symbol-transfer}
\end{equation}
with the usual outgoing, limiting-absorption, or PML regularisation near
singular frequencies. In physical space, this is a convolution operator in the
translation-invariant idealisation. This provides a second reason that
convolutional neural networks are natural approximators for homogeneous
frequency transfer: away from boundaries, the exact map has an integral-kernel
structure.

The Fourier expression also gives a warning. The multiplier can become large
near high-frequency resonant components. A learned approximation may predict
the dominant field structure well while still making small errors in
operator-sensitive components. Those errors are invisible if one looks only at
the plotted field, but they can dominate the residual.

% ============================================================
\section{Warm starts and induced residuals}
\label{sec:ch5-warm-starts}
% ============================================================

The simplest use of learned solution transfer is as a warm start. One first
solves the cheaper low-frequency problem, applies a learned transfer map
$T_{\uparrow,\theta}^{\mathrm{sol}}$, and uses the result as the initial guess
for the high-frequency solve:
\begin{equation}
    x_0
    =
    T_{\uparrow,\theta}^{\mathrm{sol}}u_L.
    \label{eq:ch5-warm-start}
\end{equation}
If $u_H$ is the exact high-frequency solution, the initial error and residual
are
\begin{equation}
    e_0
    =
    u_H-x_0,
    \qquad
    r_0
    =
    f-\mathbf{A}_H x_0
    =
    \mathbf{A}_H e_0.
    \label{eq:ch5-error-residual-relation}
\end{equation}
The residual is therefore not an independent diagnostic. It is the field error
after applying the high-frequency discrete operator.

Equation~\eqref{eq:ch5-error-residual-relation} is the first major bridge from
frequency transfer to solver performance. A field-trained model is usually
optimised to make $\|e_0\|_2/\|u_H\|_2$ small. FGMRES, however, is driven by
$\|r_0\|_2/\|f\|_2$ or by the corresponding preconditioned residual. These
quantities can behave very differently because $\mathbf{A}_H$ is indefinite
and has strongly frequency-dependent amplification.

Thus a learned field prediction can occupy three qualitatively different
regimes:
\begin{enumerate}
    \item it may be accurate as a field and also residual-compatible;
    \item it may be accurate as a field but poor as an algebraic initial state;
    \item it may improve only some components of the residual and damage
    others.
\end{enumerate}
The later results show examples of these regimes. The purpose here is to make
clear why they are mathematically possible.

For this thesis, the headline solver-facing warm-start metric is the true
relative residual
\begin{equation}
    \frac{\|f-\mathbf{A}_H x_0\|_2}{\|f\|_2}.
    \label{eq:ch5-true-initial-residual}
\end{equation}
The relative field error remains important, but it answers a different
question. Field error asks whether the predicted wavefield resembles the
target solution. The residual asks whether the prediction is a good state for
the discrete equation that the Krylov method will solve.

% ============================================================
\section{From solution transfer to preconditioning}
\label{sec:ch5-transfer-precond}
% ============================================================

A warm start uses frequency transfer once, before the Krylov iteration begins.
A preconditioner uses a transfer mechanism repeatedly inside the iteration.
This changes the mathematical task. At FGMRES iteration $j$, the input is a
high-frequency residual $r_H^{(j)}$, and the preconditioner must return a
correction $z_H^{(j)}$ such that
\begin{equation}
    \mathbf{A}_H z_H^{(j)}
    \approx
    r_H^{(j)}.
    \label{eq:ch5-preconditioner-task}
\end{equation}
The output is not a complete high-frequency solution. It is an approximate
error correction for the current residual.

This suggests a learned V-cycle structure,
\begin{equation}
    \mathbf{M}_{\theta}^{-1}
    =
    \mathcal{P}_{\theta}
    \mathbf{A}_L^{-1}
    \mathcal{R}_{\theta},
    \label{eq:ch5-learned-vcycle}
\end{equation}
where $\mathcal{R}_{\theta}$ maps a high-frequency residual to a
low-frequency right-hand side or correction representation, $\mathbf{A}_L^{-1}$
performs a low-frequency solve, and $\mathcal{P}_{\theta}$ prolongs the
low-frequency correction back to the high-frequency system. The analogy with a
geometric multigrid V-cycle is useful:
\begin{equation}
    \text{restriction}
    \;\longrightarrow\;
    \text{coarse solve}
    \;\longrightarrow\;
    \text{prolongation}.
    \label{eq:ch5-vcycle-words}
\end{equation}
The difference is that the transfer maps must respect both grid geometry and
frequency-dependent wave physics.

\begin{figure}[htbp]
    \centering
    \chfiveincludegraphics{0.82\linewidth}{figures/ch5/learned_vcycle_schematic.png}
    \caption{Learned V-cycle preconditioner inside one FGMRES iteration. The
    high-frequency residual $r_H^{(j)}$ is restricted by a learned map
    $\mathcal{R}_{\theta}$, solved with the low-frequency operator
    $\mathbf{A}_L^{-1}$, and prolonged by a learned map
    $\mathcal{P}_{\theta}$ to produce the high-frequency correction
    $z_H^{(j)}=\mathbf{M}_{\theta}^{-1}r_H^{(j)}$.}
    \label{fig:ch5-learned-vcycle-schematic}
\end{figure}

The thesis therefore distinguishes three transfer objects:
\begin{align}
    T_{\uparrow}^{\mathrm{sol}} &: u_L \mapsto u_H,
    &&\text{solution transfer}, \label{eq:ch5-three-objects-solution}\\
    R &: r_H \mapsto r_L \text{ or } e_L,
    &&\text{residual restriction}, \label{eq:ch5-three-objects-restriction}\\
    P &: e_L \mapsto e_H,
    &&\text{correction prolongation}. \label{eq:ch5-three-objects-prolongation}
\end{align}
These objects have related physical content, but they are not interchangeable.
A solution-transfer network has learned how complete solutions at two
frequencies are related on the solution manifold. A residual restriction must
map an arbitrary Krylov residual to a useful low-frequency problem. A
correction prolongation must map an approximate low-frequency error correction
to a high-frequency correction that reduces the current residual. Confusing
these roles is a direct route to a learned component that looks plausible as a
field but fails as a preconditioner.

\begin{figure}[htbp]
    \centering
    \chfiveincludegraphics{0.92\linewidth}{figures/ch5/three_transfer_objects.png}
    \caption{Three transfer objects that must be kept separate. Solution
    transfer maps one complete wavefield to another. Residual restriction maps
    a current high-frequency residual to a low-frequency representation.
    Correction prolongation maps a low-frequency correction back to the
    high-frequency system. They share frequency-transfer intuition but have
    different algebraic roles.}
    \label{fig:ch5-three-transfer-objects}
\end{figure}

% ============================================================
\section{Residual amplification and modal safety}
\label{sec:ch5-modal-safety}
% ============================================================

The cleanest way to see why field accuracy and residual accuracy differ is the
one-dimensional Dirichlet problem. Let $\mathbf{A}_H$ have eigenpairs
$(\lambda_k,v_k)$ with $\|v_k\|_2=1$. If
\begin{equation}
    e_0
    =
    u_H-x_0
    =
    \sum_k c_k(e_0)v_k,
    \label{eq:ch5-error-expansion}
\end{equation}
then
\begin{equation}
    r_0
    =
    \mathbf{A}_H e_0
    =
    \sum_k \lambda_k c_k(e_0)v_k.
    \label{eq:ch5-residual-expansion}
\end{equation}
Thus
\begin{equation}
    c_k(r_0)
    =
    \lambda_k c_k(e_0).
    \label{eq:ch5-modal-amplification}
\end{equation}
The residual coefficient is the field-error coefficient multiplied by the
corresponding eigenvalue.

For the standard second-order Dirichlet discretisation of
$-\partial_{xx}-\omega^2$, the eigenvectors are sine modes,
\begin{equation}
    v_k(j)
    =
    \sqrt{\frac{2}{N+1}}
    \sin\!\left(\frac{jk\pi}{N+1}\right),
    \qquad
    j,k=1,\dots,N,
    \label{eq:ch5-dirichlet-eigenvectors}
\end{equation}
and the eigenvalues are
\begin{equation}
    \lambda_k(\omega)
    =
    \frac{4}{h^2}
    \sin^2\!\left(\frac{k\pi}{2(N+1)}\right)
    -
    \omega^2.
    \label{eq:ch5-dirichlet-eigenvalues}
\end{equation}
This basis is used later as a diagnostic microscope. It allows one to see not
only whether a neural field has small error, but where that error sits in the
spectrum and how it is amplified into a residual.

Equation~\eqref{eq:ch5-modal-amplification} motivates the notion of modal
safety. A learned component is safe in a mode if inserting it reduces the
residual contribution in that mode relative to a trusted baseline. It is
unsafe if it improves the field visually but increases the residual
contribution. In a setting with known eigenvectors, this can be tested by a
spectral gate: accept a learned modal component only if it reduces the
corresponding residual. In higher-dimensional PML problems the same exact
modal gate is not available, but the principle remains: learned information
must be selected or trained according to residual improvement, not only field
similarity.

\begin{figure}[htbp]
    \centering
    \chfiveincludegraphics{0.86\linewidth}{figures/ch5/modal_amplification_cartoon.png}
    \caption{Residual amplification in the Dirichlet eigenbasis. A field error
    with small high-mode coefficients can still create a large residual
    because each residual coefficient is multiplied by the corresponding
    eigenvalue, $c_k(r_0)=\lambda_k c_k(e_0)$. This is the central reason that
    field accuracy alone is not a sufficient criterion for Krylov
    acceleration.}
    \label{fig:ch5-modal-amplification-cartoon}
\end{figure}

% ============================================================
\section{PML compatibility is operator compatibility}
\label{sec:ch5-pml}
% ============================================================

The preceding modal discussion is exact for a symmetric Dirichlet model
problem. The solver-facing two-dimensional system in this thesis is different:
it uses a finite-difference/PML operator on a full computational grid. The PML
region is not a visual padding region around the physical domain. It is part
of the discrete operator.

This has an important consequence for learned transfer. Suppose a network
predicts a high-frequency field that is accurate in the physical interior but
inconsistent in the absorbing layer. The interior field may look acceptable,
and the interior relative error may be small, but the full-grid residual
\begin{equation}
    r_0
    =
    b-\mathbf{A}_H x_0
    \label{eq:ch5-pml-residual}
\end{equation}
is computed with the same FD/PML operator used by the solver. If the PML
values are incompatible with the stencil, coefficient stretching, or boundary
treatment, the residual can become large even when the physical-domain plot
looks reasonable.

It is therefore not generally safe to zero the PML output of a learned field,
nor to train only on the physical interior and assume that the solver will
ignore the absorbing layer. The PML is part of the algebraic problem. A
solver-facing transfer model must either predict the full-grid PML field in an
operator-compatible way or explicitly construct a correction whose full-grid
residual is controlled.

\begin{figure}[htbp]
    \centering
    \chfiveincludegraphics{0.88\linewidth}{figures/ch5/pml_operator_compatibility.png}
    \caption{PML compatibility as an operator issue. The physical interior may
    contain the visually important wavefield, but the residual is computed by
    the full FD/PML operator. A learned prediction that mishandles the
    absorbing layer can therefore be a poor solver state even if the interior
    field appears accurate.}
    \label{fig:ch5-pml-operator-compatibility}
\end{figure}

This observation connects the one-dimensional and two-dimensional parts of the
thesis. The Dirichlet problem explains residual amplification exactly through
eigenmodes. The PML problem replaces that clean spectral picture by a
complex-valued, generally non-normal operator. The diagnostic principle,
however, is unchanged: apply the same operator that the solver applies, and
measure the residual it produces.

% ============================================================
\section{Learning objectives}
\label{sec:ch5-objectives}
% ============================================================

The most direct objective for solution transfer is the complex relative
$\ell^2$ field loss
\begin{equation}
    \mathcal{L}_{\mathrm{field}}
    =
    \frac{\|\hat{u}_H-u_H\|_2^2}
         {\|u_H\|_2^2+\varepsilon}.
    \label{eq:ch5-field-loss}
\end{equation}
The real and imaginary parts are represented as separate channels in the
network, but the norm is interpreted as the norm of the complex vector. This
objective is natural for learning $T_{\uparrow}^{\mathrm{sol}}$, and it is the
right first question: can the model predict the high-frequency wavefield from
the low-frequency wavefield?

For a preconditioner, however, the natural objective is residual-based. If a
network predicts a correction $\hat{e}$ for a residual $r$, then the solver
task is
\begin{equation}
    \mathbf{A}\hat{e}
    \approx
    r.
\end{equation}
A corresponding loss is
\begin{equation}
    \mathcal{L}_{\mathrm{res}}
    =
    \frac{\|\mathbf{A}\hat{e}-r\|_2^2}
         {\|r\|_2^2+\varepsilon}.
    \label{eq:ch5-residual-loss}
\end{equation}
This objective is more aligned with Krylov acceleration than
\eqref{eq:ch5-field-loss}. It is also more demanding, because applying
$\mathbf{A}$ exposes high-amplification components of the error.

A third option is to measure residual error in a preconditioned norm. If
$\mathbf{M}^{-1}$ is a trusted baseline preconditioner, such as complex
shifted Laplace, then
\begin{equation}
    \mathcal{L}_{\mathrm{prec}}
    =
    \frac{\|\mathbf{M}^{-1}(\mathbf{A}\hat{e}-r)\|_2^2}
         {\|\mathbf{M}^{-1}r\|_2^2+\varepsilon}
    \label{eq:ch5-preconditioned-residual-loss}
\end{equation}
asks whether the learned correction is useful in the norm seen by the
preconditioned Krylov process. This is not identical to the true residual, but
it can be a meaningful training or diagnostic quantity when the learned method
is intended to augment a baseline preconditioner.

\begin{table}[htbp]
    \centering
    \caption{Learning objectives and the mathematical question each one
    answers.}
    \label{tab:ch5-objectives}
    \begin{tabular}{@{}p{0.24\linewidth}p{0.38\linewidth}p{0.28\linewidth}@{}}
        \toprule
        Objective & Question answered & Best use \\
        \midrule
        Field loss
        & Does the predicted wavefield match the target solution?
        & Solution transfer and warm starts \\
        Residual loss
        & Does the predicted correction solve $\mathbf{A}e=r$?
        & Correction learning and learned preconditioning \\
        Preconditioned residual loss
        & Does the correction help in the baseline preconditioned norm?
        & Hybrid learned/classical methods \\
        Modal or gated diagnostics
        & In which components is the learned correction residual-safe?
        & Interpretation and robust selection \\
        \bottomrule
    \end{tabular}
\end{table}

The later chapters use these distinctions to organise the experimental
evidence. A model may be successful under a field objective but insufficient
as a solver component. Conversely, a selective correction can be valuable even
if it modifies only a small part of the field, provided that it reduces the
residual components that control convergence.

% ============================================================
\section{Architectural implications}
\label{sec:ch5-architecture}
% ============================================================

The preceding analysis suggests several architectural requirements for learned
frequency transfer.

First, the model needs a nontrivial receptive field. The phase transformation
from $\omega_L$ to $\omega_H$ depends on propagation distance and cannot be
represented by independent pointwise scaling. Convolutional architectures are
therefore natural because they approximate local integral kernels and can
aggregate neighbouring wavefield information.

Second, the representation must preserve complex phase. This thesis represents
complex fields by real and imaginary channels. An amplitude--phase
representation is possible, but it introduces discontinuities at phase wraps,
which are inconvenient for standard convolutional networks.

Third, the model needs multiscale structure. Low frequencies carry broad
propagation envelopes; high frequencies require oscillatory detail. A U-Net is
a natural architecture because the encoder aggregates large-scale context and
the decoder reconstructs fine-scale structure, while skip connections preserve
spatial detail.

Fourth, architecture alone does not guarantee solver usefulness. A network can
have enough expressive power to predict plausible fields and still produce
residual-unfriendly components. The solver-facing tests must therefore include
true residuals, preconditioned residuals, modal diagnostics where available,
and FGMRES convergence.

\begin{figure}[htbp]
    \centering
    \chfiveincludegraphics{0.90\linewidth}{figures/ch5/unet_frequency_transfer_schematic.png}
    \caption{U-Net as a frequency-transfer model. The encoder aggregates
    large-scale wavefield structure, while skip connections preserve
    oscillatory detail needed by the decoder. The same architectural family
    can be used for solution transfer or correction prediction, but the target
    and loss determine which mathematical object is learned.}
    \label{fig:ch5-unet-frequency-transfer-schematic}
\end{figure}

This thesis focuses on U-Net transfer models because they are expressive
enough to learn frequency-coherent solution maps, while still being simple
enough to diagnose. The scientific question is not whether a larger neural
operator can reduce field loss further. The question is whether learned
frequency information can be made residual-compatible with the Helmholtz
solver.

% ============================================================
\section{Summary}
\label{sec:ch5-summary}
% ============================================================

This chapter introduced the frequency transfer operator as the mathematical
link between low- and high-frequency Helmholtz fields. The Green's function
shows why such transfer is physically plausible: changing frequency modifies
phase and amplitude scale while preserving propagation geometry. Exact
solution transfer is therefore well-defined as
$\mathcal{T}_{\uparrow}^{\mathrm{sol}}=\mathcal{A}_H^{-1}\mathcal{A}_L$.

The chapter also established the central caution. A solution-transfer operator
is not automatically a preconditioner. A warm start is judged by the residual
it induces under the high-frequency operator, and an in-Krylov preconditioner
must map residuals to corrections. Because
$r_0=\mathbf{A}_H(u_H-x_0)$, small field error can become large residual error
when the error lies in operator-amplified components.

This distinction motivates the computational and experimental chapters that
follow. The computational setup specifies the discrete operators, datasets,
models, and residual metrics needed to test the idea without ambiguity. The
results then ask the solver-facing question: when does learned frequency
transfer remain useful after the high-frequency Helmholtz operator, the PML,
and the Krylov residual have all been taken seriously?

% ============================================================
% Recommended visual plan for this chapter
% ============================================================
%
% Essential figures:
% 1. figures/ch5/green_frequency_coherence.png
%    Real parts of Green's functions at several frequencies.
%    Purpose: show common propagation geometry and increasing oscillation.
%
% 2. figures/ch5/radial_green_cross_section.png
%    Radial amplitude and real-part cross sections for a doubled-frequency pair.
%    Purpose: show the 1/sqrt(2) far-field amplitude scaling and doubled phase.
%
% 3. figures/ch5/solution_transfer_operator_schematic.png
%    Diagram of f -> A_L^{-1} -> u_L -> T_up^sol -> u_H and
%    f -> A_H^{-1} -> u_H.
%    Purpose: define exact solution transfer without ambiguity.
%
% 4. figures/ch5/three_transfer_objects.png
%    Three rows: solution transfer, residual restriction, correction
%    prolongation.
%    Purpose: prevent the central conceptual confusion.
%
% 5. figures/ch5/modal_amplification_cartoon.png
%    Small field-error coefficients multiplied by lambda_k become residual
%    coefficients.
%    Purpose: foreshadow the 1D spectral diagnostics in the Results chapter.
%
% 6. figures/ch5/pml_operator_compatibility.png
%    Full-grid FD/PML operator acting on a prediction with interior and PML
%    regions marked.
%    Purpose: show that PML compatibility is residual/operator compatibility.
%
% 7. figures/ch5/learned_vcycle_schematic.png
%    r_H -> R_theta -> A_L^{-1} -> P_theta -> z_H inside FGMRES.
%    Purpose: show the learned preconditioning path.
%
% Optional figure:
% 8. figures/ch5/unet_frequency_transfer_schematic.png
%    U-Net encoder/decoder with complex channels.
%    Purpose: support the architecture section if space permits.
